\section{Recognizing Sound Shapes}\label{sec:shapes}

A spectrogram turns musical events into shapes. Read position as frequency,
horizontal extent as captured time, and color as magnitude. Start with
Average and Smooth off so their blending does not disguise the pattern.

\begin{figure}[!htp]
\centering
% Abstract time-frequency signatures; no app data or simulated UI.
\begin{tikzpicture}[x=1cm,y=1cm,font=\small]
  \foreach \x/\y/\title in {0/3.4/Steady harmonics,7/3.4/Pitch glide,0/0/Attack and decay,7/0/Filtered noise} {
    \begin{scope}[shift={(\x,\y)}]
      \fill[panelPaper] (0,0) rectangle (5.5,2.1);
      \node[anchor=west,font=\small\bfseries,text=spectreAccent] at (0,2.5) {\title};
      \draw[->,panelMuted] (0,0)--(5.7,0) node[right,font=\scriptsize] {$t$};
      \draw[->,panelMuted] (0,0)--(0,2.2) node[above,font=\scriptsize] {$f$};
    \end{scope}
  }
  \foreach \y/\a in {.45/90,.95/65,1.45/40} {
    \draw[spectreAccent!\a,line width=2pt] (.3,3.4+\y)--(5.1,3.4+\y);
  }
  \foreach \y/\a in {.2/90,.65/60,1.1/35} {
    \draw[spectreAccent!\a,line width=2pt] (7.3,3.4+\y) .. controls
      (8.5,3.4+\y) and (9,3.8+\y) .. (12.1,4.3+\y);
  }
  \fill[spectreAccent!70] (.6,.15) rectangle (.85,1.95);
  \foreach \i in {0,...,15} {
    \pgfmathtruncatemacro{\shade}{85-\i*5}
    \draw[spectreAccent!\shade,line width=2.5pt] (.85+\i*.25,.6)--(1.1+\i*.25,.6);
    \draw[spectreAccent!\shade,line width=1.5pt] (.85+\i*.19,1.2)--(1.04+\i*.19,1.2);
  }
  \fill[spectreAccent!12] (7.3,.15)--(7.3,.6).. controls (9,.6) and (9,1.8)..(12.1,1.8)--(12.1,.15)--cycle;
  \draw[spectreAccent,line width=1.6pt] (7.3,.6)..controls (9,.6) and (9,1.8)..(12.1,1.8);
  \foreach \i in {0,...,22} {
    \pgfmathsetmacro{\xx}{7.35+\i*.21}
    \draw[spectreAccent!25,line width=.6pt] (\xx,.22)--(\xx,.38);
  }
\end{tikzpicture}

\caption{Schematic sound signatures, shown in chronological order within
one sweep. Marks indicate energy; their shade is illustrative rather than a
calibrated palette. Spectre's scan line wraps when it reaches the right edge.}
\end{figure}

\begin{description}[style=nextline,leftmargin=0pt,labelwidth=0pt,labelsep=0pt]
  \item[Steady tone] A sine produces one horizontal band. A harmonic sound
  adds bands near integer multiples of its fundamental. With Linear Freq
  Scale, equally spaced harmonics are equally spaced on the frequency axis.
  \item[Pitch movement] A glide bends the bands upward or downward.
  Vibrato makes them undulate. Harmonics move together as the fundamental
  changes; frequency smoothing can blend neighboring bands.
  \item[Attack and decay] A short broadband attack produces a vertical
  mark. Resonances can continue as horizontal tails. A pitched percussion
  sound may combine a vertical attack with a falling pitch curve.
  \item[Sustained noise] Energy spreads over a frequency region instead
  of forming stable harmonic lines. A moving filter changes that region's
  boundary and may leave a bright resonant band near its cutoff.
\end{description}

\paragraph{Use shape as evidence}
A thick band can be a window's broad peak, modulation, several nearby tones,
or smoothing. It does not establish that the source is noisy. An evenly
colored region can simply be above Ceil or below Floor. Change one viewing
control at a time, then compare the sound again with minimal smoothing.
The next walkthroughs isolate these causes.
